% ============================================================
% §5 · Uso básico del software — Manual de Usuario K-QGMIP
% ============================================================

\section{Uso básico del software}

El flujo de uso de K-QGMIP es el mismo independientemente del método elegido
(KQNodes o KGeoMIP): el usuario edita un \textbf{script de entrada} dedicado a cada
método, configura los parámetros del análisis y lo ejecuta desde el terminal.
El \textit{framework} se encarga de orquestar internamente la carga de la red,
la preparación del subsistema, la ejecución del algoritmo y el guardado de resultados.

% ------------------------------------------------------------
\subsection{Flujo general de uso}
% ------------------------------------------------------------

\begin{enumerate}
  \item \textbf{Seleccionar el método}: elegir entre KQNodes o KGeoMIP según el
        tamaño del sistema y los recursos disponibles (ver Sección de Requisitos).

  \item \textbf{Abrir el script de entrada del método elegido}:
        \begin{itemize}
          \item KQNodes: \texttt{code/QNodes/exec\_kqnodes.py}
          \item KGeoMIP: \texttt{code/GeoMIP/exec\_kgeomip.py}
        \end{itemize}

  \item \textbf{Configurar los parámetros} directamente en el script
        (detallados en la subsección siguiente).

  \item \textbf{Ejecutar el script} desde la carpeta correspondiente:
\begin{lstlisting}[language=bash]
# Para KQNodes
cd code/QNodes
uv run exec_kqnodes.py

# Para KGeoMIP
cd code/GeoMIP
uv run exec_kgeomip.py
\end{lstlisting}

  \item \textbf{Revisar los resultados} generados automáticamente en
        \texttt{results/\{metodo\}/} en formato CSV y Markdown.
\end{enumerate}

% ------------------------------------------------------------
\subsection{Parámetros que debe especificar el usuario}
% ------------------------------------------------------------

Cada script de entrada expone cuatro parámetros al inicio del archivo que el usuario
debe ajustar antes de ejecutar. Tres de ellos son comunes a ambos métodos
(\texttt{ESTADO}, \texttt{K}, \texttt{MUESTRA}); el cuarto difiere según el método.

\subsubsection*{KQNodes — \texttt{code/QNodes/exec\_kqnodes.py}}

\begin{table}[H]
\centering
\begin{tabular}{|l|p{3.5cm}|p{6cm}|}
\hline
\textbf{Parámetro} & \textbf{Ejemplo} & \textbf{Descripción} \\
\hline
\texttt{ESTADO} & \texttt{"1" + "0" * 19} &
  Estado inicial $s(t)$: cadena binaria de longitud $n$.
  Define implícitamente el tamaño del sistema. \\
\hline
\texttt{K} & \texttt{5} &
  Número de partes de la k-partición. Valores válidos: \texttt{2}, \texttt{3},
  \texttt{4}, \texttt{5}. Con \texttt{2} reproduce la bipartición clásica. \\
\hline
\texttt{CRITERIO} & \texttt{"C4"} &
  Criterio de refinamiento submodular.
  \texttt{"C4"} = corte marginal mínimo (recomendado);
  \texttt{"C1"} = tamaño máximo de parte. \\
\hline
\texttt{MUESTRA} & \texttt{"A"} &
  Letra de la red de prueba (\texttt{N\{n\}\{MUESTRA\}.csv}). \\
\hline
\end{tabular}
\caption{Parámetros de configuración de KQNodes en \texttt{exec\_kqnodes.py}.}
\label{tab:params_kqnodes_5}
\end{table}

\subsubsection*{KGeoMIP — \texttt{code/GeoMIP/exec\_kgeomip.py}}

\begin{table}[H]
\centering
\begin{tabular}{|l|p{3.5cm}|p{6cm}|}
\hline
\textbf{Parámetro} & \textbf{Ejemplo} & \textbf{Descripción} \\
\hline
\texttt{ESTADO} & \texttt{"1" + "0" * 24} &
  Estado inicial $s(t)$: cadena binaria de longitud $n$.
  Define implícitamente el tamaño del sistema. \\
\hline
\texttt{K} & \texttt{5} &
  Número de partes de la k-partición. Valores válidos: \texttt{2}, \texttt{3},
  \texttt{4}, \texttt{5}. Con \texttt{2} reproduce exactamente GeoMIP. \\
\hline
\texttt{VARIANTE} & \texttt{"E4"} &
  Heurística de refinamiento geométrico.
  \texttt{"E4"} = divisivo por MinHeap con $\Delta\Phi$ (recomendado; garantiza
  regresión $k=2$ y monotonicidad);
  \texttt{"A"} = clustering aglomerativo (solo para comparación experimental). \\
\hline
\texttt{MUESTRA} & \texttt{"A"} &
  Letra de la red de prueba (\texttt{N\{n\}\{MUESTRA\}.csv}). \\
\hline
\end{tabular}
\caption{Parámetros de configuración de KGeoMIP en \texttt{exec\_kgeomip.py}.}
\label{tab:params_kgeomip_5}
\end{table}

Las máscaras binarias de \textbf{condición}, \textbf{alcance} y \textbf{mecanismo}
no se configuran manualmente en los scripts: se leen automáticamente de las
hojas del archivo Excel de pruebas (\texttt{data/DatosPruebas2026\_1.xlsx}),
donde están expresadas como letras (\texttt{"ABCDE"}) y el \textit{framework}
las convierte a cadenas binarias en tiempo de ejecución.

% ------------------------------------------------------------
\subsection{Formas de proporcionar los parámetros}
% ------------------------------------------------------------

K-QGMIP admite dos formas de entrada de datos:

\begin{itemize}
  \item \textbf{Directamente en el script} (análisis individual): el usuario edita
        \texttt{ESTADO}, \texttt{K}, \texttt{CRITERIO} y \texttt{MUESTRA} en el script
        de entrada y ejecuta una sola corrida. Útil para explorar un sistema específico
        o depurar resultados.

  \item \textbf{Mediante archivo Excel} (análisis en lote): el \textit{framework}
        lee todas las combinaciones de alcance y mecanismo definidas en
        \texttt{DatosPruebas2026\_1.xlsx} para el tamaño $n$ correspondiente, y ejecuta
        el método sobre cada una. El usuario solo configura \texttt{ESTADO}, \texttt{K},
        \texttt{CRITERIO} y \texttt{MUESTRA}; el resto lo gestiona el \textit{framework}
        automáticamente. Esta es la modalidad de ejecución por defecto al correr los
        scripts de entrada.
\end{itemize}
